  % ---------------------------------------------------------------------------
  \section{Introduction}
  \label{sec:introduction}
  % ---------------------------------------------------------------------------

  Consider a controller agent that receives a worker's private report and then
  chooses between a calculator and a search API. The report can pass a
  transcript-level audit---it need not leak the API label verbatim---while a
  counterfactual swap of the worker report changes the controller's realized
  action distribution by $\MultiAgentCounterfactualBits$ bits. This is the
  audit failure mode: the visible transcript can look adequate while an
  unobserved routing or evidence state remains both unrecoverable and
  decision-relevant. A weaker version appears in ReAct scratchpads: controlled
  replay changes the soft tool-token distribution without changing argmax tool
  counts. That is useful as a policy-readout diagnostic, but it is not the
  headline behavioral result. Identifying which hidden channels are merely
  present, and which ones actually drive decisions, is the audit task.
  How much internal state is unrecoverable from the trace, and how much of that
  state drives the next action, are distinct quantities. Treating them as a
  single ``hiddenness'' judgment produces something neither computable nor
  diagnostic: a system with unrecoverable but dormant hidden state is different
  from one whose unrecoverable state changes the action distribution, even when
  the visible transcript looks equally complete.

  We propose a dual-certificate protocol that answers each question separately.
  The non-recoverability certificate
  $\varepsilon_\mathrm{unrec}^\mathrm{LB}(t)$ lower-bounds residual hidden-state
  uncertainty $H(S_t \mid \tilde T_t)$ by measuring finite witnesses of the
  hidden state or by evaluating closed-form finite-dimensional witness models.
  The action-relevance certificate
  $\delta_\mathrm{act}^\mathrm{LB}(t)$ lower-bounds
  $I(S_t;A_t\mid\tilde T_t)$ from admissible replay, intervention, or proxy
  measurements using conditional data processing. A capacity map with min-cut
  value $\varepsilon_\text{state}^\text{UB}$ is still reported when
  architecture access is available, but its role is an upper envelope on
  residual latent-channel capacity, not the main security claim. The central audit
  question is whether
  \emph{both}
  \[
    H(S_t\mid\tilde T_t)\geq\varepsilon_\mathrm{unrec}^\mathrm{LB}(t),
    \qquad
    I(S_t;A_t\mid\tilde T_t)\geq\delta_\mathrm{act}^\mathrm{LB}(t)
  \]
  hold with nontrivial lower bounds.
  In real systems where no direct finite witness of $S_t$ is exposed, a
  nonzero $\delta_\mathrm{act}^\mathrm{LB}$ still certifies
  non-recoverability through the elementary bound
  $I(S_t;A_t\mid\tilde T_t)\leq H(S_t\mid\tilde T_t)$. We call this an
  \emph{action-mediated} non-recoverability certificate. It is weaker than an
  independent state-only witness, but it is exactly the gap an audit cares
  about: hidden state omitted from the transcript changes the action.

  The dynamic axis can also be indexed. Instead of asking only whether residual
  state matters, an auditor can ask \emph{where} it matters: which module,
  denoising step, or communication edge carries decision relevance. This turns
  the certificate from a scalar diagnostic into an activation profile over the
  hidden-state geometry of the deployment.

  Conditioning on $\tilde T_t$ is part of the protection, not a nuisance term.
  If a model writes a steganographic or compressed rationale into the recorded
  trace, that information is in $\tilde T_t$ and is conditioned away by $I(S_t;
  A_t \mid \tilde T_t)$. The dynamic certificate measures residual decision
  relevance beyond the logged trace; it does not count hiddenness that has
  already been made visible to the audit record.

  \textbf{Contributions.} First, we separate transcript recoverability from
  residual action relevance as two distinct audit questions. Second, we give
  dual lower certificates
  $(\varepsilon_\mathrm{unrec}^\mathrm{LB},
  \delta_\mathrm{act}^\mathrm{LB})$ for those questions. Third, under explicit
  audit-abstraction assumptions, we prove certificate validity: finite
  hidden-state witnesses lower-bound $H(S_t\mid\tilde T_t)$, and admissible
  probes lower-bound $I(S_t;A_t\mid\tilde T_t)$. Fourth, we give closed-form or
  computable audit instances for diffusion trajectories, retrieval-style hidden
  source selection, and nested agent architectures. An auxiliary deployment
  capacity map with min-cut value $\varepsilon_\text{state}^\text{UB}$
  identifies where hidden capacity can reside and supports logging ablations,
  but it is not a lower certificate.

  We instantiate the protocol on three audit geometries: a ReAct scratchpad
  module, a diffusion-LM denoising trajectory, and a multi-agent private-report
  edge. These real-system diagnostics report action-mediated lower certificates
  alongside the capacity map with min-cut value
  $\varepsilon_\text{state}^\text{UB}$.
  Appendix calibrations cover a synthetic ground-truth setting where an
  independent hidden-state witness makes
  $\varepsilon_\mathrm{unrec}^\mathrm{LB}$ explicit and the capacity-map
  min-cut value $\varepsilon_\text{state}^\text{UB}$ matches true hidden-state
  entropy.
